% Ad Atticum 2.3 (ad-atticum-02-03) --- English translation
% Translated by Claude (Anthropic) from the Latin (Perseus).
% Style guide: STYLE.md.

\ciceroLetterOpener{Cicero to Atticus, greetings}

\ciceroSection{1} First, I take it, the good news \textit{[Greek: euangelia]}: Valerius has been acquitted, with Hortensius defending. That trial was thought to have been donated to the son of Aulus, and I suspect, as you write, that Epicrates has been wanton; for indeed his stockings and white-bound headbands did not please me. What it is we shall know when you come.

\ciceroSection{2} Where you find fault with the narrowness of the windows: know that you find fault with the \textit{Cyropaedia} \textit{[Greek: Kyrou paideian]}. For when I said the same thing myself, Cyrus said that with broad apertures the views \textit{[Greek: diaphaseis]} of the gardens were not so sweet. For let A \textit{[Greek: \=e a]} be the eye, B \textit{[Greek: to de hor\=omenon]} the object seen, the rays \textit{[Greek: aktines de]} D and E. You see the rest. For if we saw by emanations \textit{[Greek: kat' eid\=ol\=on]}, the emanations would labour greatly in narrow places; as it is, the outpouring \textit{[Greek: ekchysis]} of rays comes off prettily. If you find fault with anything else, you will not bear me silent unless it shall be of such a kind that it can be put right without expense.

\ciceroSection{3} I come now to the month of January and to our basis \textit{[Greek: hypostasin]} and policy \textit{[Greek: politeian]}, in which I argue Socratically on each side \textit{[Greek: S\=okratik\=os eis hekateron]}, but yet at the end (as those Socratics were wont to do) the pleasing one \textit{[Greek: t\=en areskousan]}. The matter is one of great counsel. For either we must bravely resist the agrarian law (in which there is some struggle, but it is full of praise); or we must keep quiet, which is much like going off to Solonium or Antium; or we must even help, which Caesar, they say, expects of me, with no doubt of it. For Cornelius came to me --- I mean Balbus, Caesar's intimate. He affirmed that Caesar would in all things use my counsel and Pompey's, and would give pains to join Crassus with Pompey.

\ciceroSection{4} What lies on this side: a supreme connection with Pompey, also if I please with Caesar; a return to favour with my enemies; peace with the multitude; the leisure of old age. But I am moved by my own concluding sentence in book three: ``Meanwhile, the courses which from the first part of your youth, and which indeed as consul, by virtue and spirit you sought ---/ keep these and increase them, the fame and the praises of the good.'' When this was set down for me by Calliope herself, in that book in which many things are written aristocratically \textit{[Greek: aristokratik\=os]}, I do not think I should doubt that it always seems to us: ``one omen is best, to fight for one's country'' \textit{[Greek: heis oi\=onos aristos amynesthai peri patr\=es]} (Iliad 12.243). But let us save these things for our walks at the Compitalia. Be sure on the day before the Compitalia. I shall give orders for the bath to be heated. Terentia is asking Pomponia; we shall add the mother. Bring me the book of Theophrastus \textit{[Greek: Theophrastou peri]} from your brother Quintus's books.
